% Abstract for the chapter "Parameterized Verification".
\begin{abstract}
We consider the model checking problem for \emph{parameterized systems}:
systems that consist of an arbitrary number of components. The task is to show
correctness regardless of the number of components. Examples of parameterized
systems include mutual exclusion algorithms, cache coherence protocols,
network protocols, and distributed algorithms. The problem is undecidable in
general, but a wide range of techniques has been developed that achieve
decidability or effective automation for practically important classes of
systems.

We present three techniques in detail. Regular model checking represents
configurations as words and transition relations as finite-state transducers,
and we describe an acceleration technique that computes the transitive closure
of a transducer by quotienting. Monotonic abstraction over-approximates
systems with global conditions by well-structured transition systems, so that
safety can be checked by backward reachability analysis on upward-closed sets.
View abstraction inspects only small instances of the system and recombines
their views to obtain an over-approximation of all instances, detecting a
cutoff during the analysis; it is complete for a large class of well-structured
systems. We then relate these techniques to other approaches, including cutoff
theorems, abstraction and compositional reasoning, and deductive and SMT-based
methods, and to other chapters of this handbook, and we discuss applications to
programs running under weak memory models and to parameterized quantum
circuits.
\end{abstract}
